\SourceLeaf{152}
Anmerkung in den A. N. Nr. 110, S. 230\HistoricalNote{[*)]}{[*)}{Vgl. Band VI, S. 457.]}, im Sinn —). Er hat jedoch den zu erwartenden mittleren Fehler der Ergebnisse nicht mitberechnet. Außerdem erweckt er durch einen Umstand, den Herr \textsc{Ivory} zu Recht kritisiert, den Anschein einer viel größeren Genauigkeit der Endergebnisse, als sie wirklich haben. Er stellt nämlich die verschiedenen Ergebnisse aus verschiedenen Kombinationen zusammen, wobei seine Beobachtungen in der Nähe des Äquators stets einen überwiegenden Einfluss behalten müssen. Tatsächlich ähnelt das in gewissem Maß, wenn auch bei Weitem nicht im gleichen Umfang, dem Vorgehen des verstorbenen \textsc{Ende}, der zuweilen Polhöhen nach der Methode von \textsc{Douwes} bestimmte. Dabei verband er immer \textls{dieselbe} Höhe um die Mittagszeit mit \textls{verschiedenen} weit davon entfernten Höhen und erhielt dann Ergebnisse, die innerhalb von $0'',01$ übereinstimmten.

Auch diesen Aufsatz von \textsc{Ivory} kann man ebenso schwer Punkt für Punkt widerlegen, da ihm die logische Ordnung ebenso sehr fehlt. Ich finde darin nicht den geringsten Grund gegen die Zweckmäßigkeit der Methode der kleinsten Quadrate. Was soll man aber zu dem Verfahren sagen, das \textsc{Ivory} stattdessen anwenden will, indem er die kürzeste Pendellänge mit vielen in großer Breite kombiniert? Das ist eines verständigen Mathematikers ganz unwürdig und eigentlich das Verfahren von \textsc{Ende}\HistoricalNote{[**)]}{[**)}{Die vorstehende Kritik bezieht sich auf folgende Arbeiten von J. \textsc{Ivory}:
\par\hangindent1.5em 1. On the method of least squares. Phil. Mag. (Tilloch), \textsc{lxviii}, 1826, S. 161—165.
\par\hangindent1.5em 2. On the Ellipticity of the Earth, as deduced from Experiments made with the Pendulum. Phil. Mag. (Tilloch), \textsc{lxviii}, 1826, S. 3—10, 92—101, 241—245, 246—251, 321—326, 350—353.]}. . . . . . . .

\Needspace{6\baselineskip}
\begin{center}\textsc{Gauss} an \textsc{Encke}. Göttingen, 23. August 1831.\end{center}

. . . . . . . . . . . . . .

Mit Interesse habe ich Ihrem Brief entnommen, welchen Weg Sie eingeschlagen haben, um die Verwendung des arithmetischen Mittels zu rechtfertigen. Ich finde diesen Weg sehr überzeugend, sofern die Frage, was zu tun ist, ganz unabhängig von allen Überlegungen der\space%
